\documentclass[a4paper]{article}


\usepackage[student]{Lecture-TD}

%Changer l'interligne par défaut.
%\renewcommand{\baselinestretch}{1.5}

\lhead{Travaux dirigés - Thermochimie}% des accumulateurs
\lfoot{IUT d'Orsay - BUT 2A - 2024/2025} \cfoot{
\includegraphics[height=0.5cm]{cc-by-nc-sa.png}}  \rfoot{Jean-François Olivieri \quad \thepage}

\begin{document}

% Define base path variable  
\def \BasePath{/home/jeff/Documents/IUTOrsay/1-Cours/3p06-Cours_Batterie/01-Polycopies/02-TD/}

\begin{center}
\begin{minipage}{0.7\textwidth}
\centering
\Large \textbf{TD n$^\circ ~ 3$} \\
Utilisation de la thermochimie dans la prédiction des performances des accumulateurs
\end{minipage}
\end{center}

\objectives{
\begin{itemize}
\item Connaître la réactivité du lithium métallique au contact de l'atmosphère ou de solutions aqueuses ;
\item Justifier thermodynamiquement le potentiel très réducteur du lithium et mettre en perspective avec son utilisation dans les accumulateurs ;
\item Montrer qu'une réaction est spontanée/exergonique (resp. endergonique) en décharge (resp.charge) ;
\item \'{E}tablir la f.e.m. d'une cellule à l'aide de la relation de Nernst ou des grandeurs thermodynamiques tabulées (enthalpies de formation, entropies molaires$\ldots$) ;
\item Savoir exprimer et exploiter l'influence de la température sur la f.e.m. ;
\item Utiliser les grandeurs associées au travail électrique, à la puissance électrique, à la capacité et estimer un rendement défini par l'énoncé ; 
\item Savoir établir des relations simples entre grandeurs électriques : charge/courant et avancement d'une réaction.
\end{itemize}
}

\cteacher{Dans cette partie, on étudie quelques réactivités typiques du lithium. Un petit point sur la haute réactivité du lithium aux gaz atmosphériques (\ce{O2}, \ce{N2}) et à l'eau est à faire au début de cette partie en s'appuyant sur les exercices \ref{subsection:reac_Li_m} et \ref{subsubsection:thermochimie_du_lithium}.  Si possible, il faudrait mettre en perspective avec les synthèses réalisées en TP où nous n'aurons pas besoin de boîte à gants puisque nous serons emmené à utilisé du lithium uniquement sous sa forme cationique. \\
Vous pouvez utiliser la vidéo suivante : \url{https://www.youtube.com/watch?v=5mvWQdad31o&ab_channel=mopatin} ou \url{http://www.chim.lu/ch2343.php} pour montrer que les processus de réaction du lithium avec \ce{N2}, \ce{O2} et \ce{H2O} sont assez fulgurant et rapide. \\
Vous pouvez faire un petit point en expliquant que les propriétés du lithium dans la famille des alcalins ne sont pas dû à sa position dans le tableau périodique, mais bien à l'influence de son environnement chimique sur ses propriétés rédox. L'Ex. \ref{subsection:born-haber} reprend un passage du cours où on montre l'importance de la stabilisation de \ce{Li+} en milieu aqueux dans la valeur du potentiel rédox standard. Vous pouvez donc ouvrir sur l'objectif du module : l'importance de la substitution pour moduler la réactivité des matériaux lithiés. \\
Pour les exercices prioritaire, vous pouvez traiter en priorité les Ex. \ref{subsubsection:thermochimie_du_lithium} qui traite thermodynamiquement la réactivité du lithium dans l'eau, et l'Ex. \ref{subsubsection:chimie_des_batteries_lithium-ion} qui permet d'étudier deux matériaux d'insertion qu'ils rencontreront beaucoup par la suite : la graphite et le dioxyde de cobalt. Cet exercice aborde à peu près toute la gamme de questions que je pourrai mettre en partiel sur cette partie.
}

\section{Réactivité du lithium métallique}

\input{\BasePath Ex05}

\input{\BasePath Ex04}

\input{\BasePath Ex38}
\section{Performances du lithium dans les accumulateurs}

\input{\BasePath Ex22}

\input{\BasePath Ex09}

%\subsection{Comparaison des performances de quelques technologies de batterie}

\setcounter{equation}{0}
\setcounter{Q}{0}

Les performances d'une batterie sont jugées en premier lieu à la densité massique d'énergie. En effet, on cherche en général à stocker un maximum d'énergie dans un dispositif léger qui prend le minimum de place.%avec un minimum d'encombrement.
\begin{center}
\begin{adjustbox}{width=0.95\textwidth}
\begin{tabular}{cccccc}
    \toprule
     \textbf{Technologie} & \textbf{Année de développement} & \textbf{\'{E}nergie massique} & \textbf{\'{E}nergie volumique} & \textbf{Tension} & \textbf{Cyclabilité}\\
     Unité & & \si{\watt\hour\per\kilo\gram} & \si{\watt\hour\per\liter} & \si{\volt}/au couple \ce{Li^+}/\ce{Li} & $\emptyset$ \\
     \midrule
     Plomb & 1859 & 35 & 70 & 2.0 & 200-250\\
     Nickel-Cadmium & 1900 & 60 & 100 & 1.2 & 400-500\\
     Nickel-Métal-Hydrure & 1975 & 80 & 240 & 1.2 & 400-500\\
     Lithium-Ion & 1978 & 250 & 400 & 3.6 & > 1000\\
     \bottomrule
\end{tabular}   
\end{adjustbox}
\captionof{table}{\small \'{E}volution de l'ordre de grandeur de l'énergie massique, volumique, de la tension ainsi que de la cyclabilité d'une batterie en fonction de la technologie.}
\label{tab:battery_time}
\end{center}

\Q Définir la densité massique d'énergie et la densité volumique d'énergie. Définir le wattheure et rappeler sa relation aux unités du système international. 
\solution{La densité massique/volumique d'énergie électrique délivrable par une batterie correspond au produit : 
\begin{align*}
\text{tension de la cellule} \times \text{capacité massique/volumique}
\end{align*}
Cette grandeur permet donc de condenser deux informations cruciales : le nombre d'électrons échangeables et le "potentiel" électrique de chacun de ces électrons ! \\
Le joule \si{\joule} est l'unité de l'énergie dans le système international. Le watt-heure représente l'énergie délivrée par un dispositif électrique fournissant une puissance électrique de \SI{1}{\watt} pendant \SI{1}{\hour}. Or, il y a \SI{3600}{\second} dans \SI{1}{\hour}, d'où  
\begin{align*}
\SI{1}{\watt\hour} = \SI{3600}{\joule}
\end{align*}}

\Q Déterminer la capacité spécifique massique et volumique moyenne de chacune des technologies.
\solution{On note $CS_m$ (resp. $CS_V$) les capacités spécifiques massiques (resp. volumique) de chacune de ces technologies.
\begin{center}
\begin{adjustbox}{width=0.95\textwidth}
\begin{tabular}{cccc}
    \toprule
     \textbf{Technologie} & \textbf{Année de développement} & \textbf{Capacité massique} & \textbf{Capacité volumique} \\
     Unité & & \si{\ampere\hour\per\kilo\gram} & \si{\ampere\hour\per\liter} \\
     \midrule
     Plomb & 1859 & 17.5 & 35 \\
     Nickel-Cadmium & 1900 & 50 & 84\\
     Nickel-Métal-Hydrure & 1975 & 67 & 200\\
     Lithium-Ion & 1978 & 70 & 111\\
     \bottomrule
\end{tabular}   
\end{adjustbox}
\end{center}}

\Q Commenter les valeurs de la Tab.\ref{tab:battery_time}. Quel est l'intérêt de la technologie Li-ion par rapport à d'autres technologies ?
\solution{On peut soulever les avantages de la batterie Li : 
\begin{itemize}
\item l'énergie massique la plus grande ;
\item l'énergie volumique la plus grande ;
\item la meilleure f.e.m. ;
\item la meilleure cyclabilité.
\end{itemize}
Les batteries lithium sont donc parfaitement adaptées aux dispositifs portatifs nécessitant de grandes quantités d'énergies (\emph{par ex.} les portables).  \\
Il est à noter toutefois que la technologie Ni-M-H présente des capacités spécifiques volumiques et massiques équivalentes, voire plus intéressantes, que la technologie lithium ce qui peut offrir des performances intéressantes pour les applications nécessitant de faibles capacités en énergie (\emph{par ex.} les terminaux de paiement bancaire, les technologies NFC, les accessoires d'Airsoft, certains capteurs). 
}

\Q Rappeler ce qu'est la qualité de l'énergie stockée et expliquer l'intérêt des batteries Li. Expliquez-en l'origine.

\solution{La qualité d'énergie est une façon d'exprimer l'origine de la puissance électrique qui est soit potentielle (venant de la f.e.m.), soit cinétique (venant du courant électrique). En effet, dans un modèle de Thévenin simple (cf cours), la puissance électrique est donnée par : 
\begin{align*}
P = e I - R I^2
\end{align*}
Un dispositif présentant une grande tension à vide $e$ présente l'avantage de délivrer moins de courant pour une même puissance électrique, ce qui a pour conséquence :
\begin{itemize}
\item de diminuer les déperditions d'énergie, sous forme thermique, dû aux phénomènes résistifs (terme en $-R I^2$) ;
\item d'augmenter la durée de décharge (c.-à.-d. que la batterie se décharge plus lentement).
\end{itemize}}

\Q Mettre en relation ces données avec la place de l'élément Li au sein du tableau périodique.
\solution{Le lithium se situe en haut à gauche du tableau périodique ($3^e$ élément du tableau). C'est un élément métallique, parmi les plus électropositifs (explique la f.e.m.), de faible poids moléculaire (explique les capacités massiques) et de faible nombre quantique principal (soit un élément de petite taille, ce qui explique les capacités volumiques).}

%\subsection{Caractéristiques d'une pile commerciale}
%Panasonic

\setcounter{equation}{0}
\setcounter{Q}{0}


Panasonic Energy{\copyright} commercialise l'accumulateur ML2020, appelé \emph{Coin-type Manganese Rechargeable Lithium Batteries}, dont les spécifications techniques suivantes sont fournies : 
\begin{center}
\begin{minipage}{0.70\textwidth}
\begin{tabular}{cc}
\toprule
Spécifications & Valeurs prises \\
\midrule
Tension de charge & \SI{2.8}{\volt}$\sim$\SI{3.2}{\volt} \\
Tension nominale & \SI{3.0}{\volt} \\
Capacité nominale & \SI{45}{\milli\ampere\hour} \\
Courant de décharge & \SI{0.12}{\milli\ampere} \\
Dimensions & D : \SI{20.0}{\milli\meter} $\times$ H : \SI{2.0}{\milli\meter} \\
Poids & \SI{2.20}{\gram} \\
Nombre de cycles de vie avant $SoH = 10 \si{\percent}$ & $\approx 1000$ \\
Température de fonctionnement & -20 à \SI{60}{\celsius} \\ 
\bottomrule
\end{tabular}
\end{minipage}\begin{minipage}{0.28\textwidth}
\centering
\includegraphics[width=0.8\textwidth]{Ex24/Fig04}
\end{minipage}
\captionof{table}{Spécifications techniques de l'accumulateur ML2020 de Panasonic{\copyright} [la capacité nominale correspond à la capacité de décharge obtenue lorsque la tension vaut \SI{2.0}{\volt}.]}
\end{center}

D'après le commercial, cette batterie présente une capacité élevée ce qui en fait un dispositif idéal destiné aux clés, télécommandes ou autres objet nécessitant de l'énergie sur le long terme.

L'objectif est d'interpréter quelques expériences ayant conduit à ces spécifications techniques.

\paragraph*{Analyse de la décharge}
Lors de cycles de décharge, un expérimentateur a suivi l'évolution de la tension au sein d'un accumulateur ML2020 neuf dans les conditions expérimentales suivantes : 
\begin{itemize}
\item Charge : un générateur de tension de 3.1 \si{\volt}, une résistance de \SI{160}{\ohm} pendant 60 heures ;
\item Décharge : une résistance de \SI{20}{\kilo\ohm} (équivaut à un courant de décharge de \SI{120}{\micro\ampere}) jusqu'à ce que la tension aux bornes de l'accumulateur vale \SI{2}{\volt}.
\end{itemize}
L'ensemble a été réalisé à 4 températures différentes.

\begin{minipage}{0.48\textwidth}
\includegraphics[width=\textwidth]{Ex24/Fig02}
\end{minipage}\hfill\begin{minipage}{0.48\textwidth}
\captionof{figure}{Courbe de décharge à 4 températures différentes : -20, 0, 20 et \SI{60}{\celsius}.}
\end{minipage}

\Q Représenter schématiquement l'état de charge de la batterie lorsque la capacité est 0, 22 et 44.
\solution{L'état de charge représente le rapport entre la quantité de charges stockées sur la quantité de charges maximale dans la batterie, on peut donc représenter les symboles suivants : 
\begin{center}
\begin{tabular}{cc}
\toprule
Capacité du graphe [\si{\milli\ampere\hour}] & Schéma \\
\midrule
0 & \faBattery[4] \\
22 & \faBattery[2] \\
44 & \faBattery[0] \\
\bottomrule
\end{tabular}\end{center}
}

\Q Comment évolue la quantité maximale de charge stockable dans la batterie lorsque la température augmente. Pourquoi ?
\solution{On lit les valeurs suivantes : 
\begin{center}
\begin{tabular}{ccccc}
\toprule
Température [\si{\celsius}] & -20 & 0 & 20 & 60 \\
\midrule
Capacité [\si{\milli\ampere\hour}] & 40 & 42 & 44 & 44 \\ 
\bottomrule
\end{tabular}
\end{center}
La cinétique chimique est à l'origine de la diminution de la température lors d'un processus rédox de charge ou décharge.
On rappelle que la loi d'Arrhénius s'écrit : 
\begin{align*}
k(T) = A \exp \left( - \frac{E_A}{RT} \right)
\end{align*}
avec $A$ le facteur préexponentiel et $E_A$ l'énergie d'activation (grandeur strictement positive).
Par conséquent, une augmentation en température favorise cinétiquement le processus rédox et permet une charge plus rapide, donc une meilleure capacité.\\
\faWarning ~ Il est important de préciser trois choses supplémentaires aux étudiants : 
\begin{enumerate}
\item lors d'un refroidissement, la résistance interne augmente à cause de la perte de "mobilité" (notion non abordée en cours) des ions, illustrable à l'aide de la formule de Thévenin ;
\item si l'élevation de la température est trop grande, d'autres réactions chimiques indésirables sont aussi favorisés cinétiquement (dégradation de l'électrolyte, réaction rédox des métaux de l'électrode$\ldots$) ce qui peut causer une diminution de la capacité à plus haute température. En effet, une partie du courant de charge servira à alimenter ces processus rédox indésirables ; 
\item lors d'un échauffement, la durée de vie de la batterie diminue (c.-à-d. $SoH$ diminue)
\end{enumerate} 
} 
%La batterie décharge moins de charge à basse température, probablement à cause de la perte de mobilité des ions lithium.}

\Q Lors de la charge et de la décharge, il a été observé que la tension moyenne de charge ($\sim \SI{3.0}{\volt}$) est plus grande que la tension moyenne de décharge. Quel est l'origine de ce processus ? Une explication énergétique est attendue.
%La tension nominale, c'est-à-dire lorsque $SoC=100 ~ \si{\percent}$, correspond approximativement à la valeur qu'on aurait lors d'une charge. Qu'observe-t-on lors de la décharge ? Quel peut-être l'explication ?
\solution{La tension de décharge est plus faible que la tension de charge. En effet, lors de la charge, l'expérimentateur est emmené à devoir fournir plus d'énergie que l'énergie chimique stockable pour permettre au lithium de migrer au sein de l'électrolyte, lui permettre de changer de phase$\ldots$
Lors de la décharge, une partie de l'énergie chimique libérée est perdue pour les mêmes raisons, ce qui résulte en une tension de décharge plus faible que la tension de charge.}

\Q \`{A} partir de quel état de charge $SoC$, la tension de décharge diminue-t-elle ? Pourquoi ? Commenter l'influence de la température.%Commenter l'évolution de la tension de décharge avec la température.
\solution{Tant que la batterie est chargée $SoC > 70 ~ \si{\percent}$, la tension de la pile est relativement constante avec la température. Ce qui traduit un coefficient thermique en \si{\volt\per\kelvin} très faible. On rappelle que le coefficient thermique s'exprime comme :
\begin{align*}
e = - \frac{\Delta_r H}{mF} + \frac{\Delta_r S}{mF} T \iff \left( \frac{\partial e}{\partial T} \right)_{P, \xi} = \frac{\Delta_r S}{mF}
\end{align*}
Pour $SoC < 30 ~ \si{\percent}$, la chute de potentiel est dû à l'augmentation de la résistance interne/surtension de l'accumulateur. Microscopiquement, il devient de plus ne plus difficile pour les ions \ce{Li^+} de quitter l'anode et de s'insérer dans la cathode.
}


\paragraph*{\'{E}volution de l'état de santé de l'accumulateur en fonction du nombre de cycles de charge/décharge}

Dans les mêmes conditions que précédemment, l'accumulateur a été soumis à 10000 cycles de charge/décharge.

\begin{minipage}{0.48\textwidth}
\includegraphics[width=\textwidth]{Ex24/Fig01}
\end{minipage}\hfill\begin{minipage}{0.48\textwidth}
\captionof{figure}{\'{E}volution de l'état de santé d'une batterie après $N$ cycles de charge/décharge.}
\end{minipage}

\Q Sachant qu'un condensateur peut faire des centaines de milliers, voire des millions d cycles, comparer cette technologie à celle de la batterie lithium. 
\solution{
Au bout de 100 cycles, cette batterie lithium a perdu \SI{50}{\percent} de ses capacités de stockage. On est très loin des performances de stockage d'un condensateur.
Cette différence est assez simple à comprendre. Le stockage d'énergie d'un condensateur est d'origine physique (confinement d'un champ électrique entre deux plaques chargées), ce qui n'induit pas de modifications chimiques pourvant être des sources microscopiques d'irréversibilité. Au contraire, les batteries sont basées sur le stockage chimique qui s'accompagne nécessairement d'une multitude de sources d'irréversibilité.
}

\Q Vérifier les spécifications $SoH$ du constructeur.
\solution{On observe qu'après 1000 cycles de charge/décharge, l'état de santé de la pile est de \SI{10}{\percent}. Ceci est en accord avec les informations du constructeur.}

Le constructeur garantit les spécifications fournies pour une tension de charge inférieure à \SI{3.3}{\volt}.

\Q Quel serait l'allure de ce graphique si la charge de l'accumulateur avait lieu à une plus grande tension ? Pourquoi ?
\solution{On doit s'attendre à avoir un déclin de l'état de santé de la batterie bien plus rapide. On peut proposer la courbe suivante : 
\begin{center}
\includegraphics[width=0.6\textwidth]{Ex24/Corr01}
\end{center}
On s'attend à avoir une dégradation des matériaux sous l'effet d'une charge $I$ trop importante. \\
\textbf{Nota Bene : } Vous pouvez commenter l'effet de température sur l'évolution du $SoH$ avec le nombre de cycle.
}

\begin{comment}
\paragraph*{\'{E}volution de la capacité en charge continue à \SI{60}{\celsius}}

La capacité 
\begin{itemize}
\item Charge : un générateur de tension de 3.1 à 3.3 \si{volt}, 
\item Décharge : une résistance de \SI{1.5}{\kilo\ohm} jusqu'à ce que la tension aux bornes de l'accumulateur vale \SI{2}{\volt}.
\end{itemize}


\begin{minipage}{0.48\textwidth}
\includegraphics[width=\textwidth]{Ex24/Fig03}
\end{minipage}\hfill\begin{minipage}{0.48\textwidth}
\captionof{figure}{Charge continue à \SI{60}{\celsius} entre 3.1 et \SI{3.3}{\volt}.}
\end{minipage}
\end{comment}

\end{document}
